% META: {"id": "1SPE-TRIGO-EX-035", "chapitre": "1SPE-TRIGONOMETRIE", "type_objet": "exercice", "capacites": ["1SPE-TRIGONOMETRIE-C4"], "capacites_codes": ["C4"], "methodes": ["M4"], "parcours": 2, "competences": ["calculer", "raisonner"], "duree_min": 12, "mode_creation": "ex_nihilo", "sources_inspiration": [], "parametres_sympy": {}, "fichier_tex": "chapitres/1SPE-TRIGONOMETRIE/exercices/1SPE-TRIGO-EX-035.tex", "corrige_tex": "chapitres/1SPE-TRIGONOMETRIE/corriges/1SPE-TRIGO-CO-035.tex", "status": "generated"}
% BEGIN-VERIFY
% from sympy import *
% x = symbols('x', real=True)
% # cos(2x) = 1/2 => 2x = pi/3 + 2k*pi or 2x = -pi/3 + 2k*pi
% # => x = pi/6 + k*pi or x = -pi/6 + k*pi
% assert cos(pi/3) == Rational(1,2)
% # On [0;2pi]: pi/6, 5pi/6, 7pi/6, 11pi/6... wait
% # x=pi/6, x=pi/6+pi=7pi/6, x=-pi/6+pi=5pi/6, x=-pi/6+2pi=11pi/6
% assert cos(2*pi/6) == cos(pi/3)
% assert cos(2*Rational(5,6)*pi) == cos(5*pi/3)
% END-VERIFY
\begin{exercice}{1SPE-TRIGO-EX-035}{2}{12}
Resoudre sur $[0\,;\,2\pi]$ les equations suivantes.
\begin{enumerate}
  \item $\cos(2x) = \dfrac{1}{2}$
  \item $\sin(2x) = -\dfrac{\sqrt{3}}{2}$
\end{enumerate}
\end{exercice}
